\documentclass{article}
\usepackage[a4paper, margin=2cm]{geometry} 
\usepackage[T1]{fontenc} %lädt Deutsche schriftzeichen
\usepackage{lmodern} % Font
\usepackage{amssymb}%für Mathe
\usepackage{setspace} %für Zeilen Abstand 
\usepackage[none]{hyphenat} % Silbentrennung deaktivieren 
\usepackage{fancyhdr}%Kopf und Fußzeile: https://de.overleaf.com/learn/latex/Headers_and_footers
\usepackage{graphicx}%Zum Bilder einfügen
\usepackage{xcolor}
\usepackage{wrapfig}
\usepackage[document]{ragged2e}%für Blocksatz
\usepackage[ddmmyyyy]{datetime}%zum bearbeiten vom Datungs Format
\usepackage{eurosym}%für \Euro / \euro
\usepackage{hyperref}
\usepackage{listings}


\newcommand{\authoren}{Niklas Hainz}
\newcommand{\Thema}{}
\newcommand{\Fach}{}
\newcommand{\Zeilenabstand}{1.2}
\newcommand{\Inhaltsverzeichnisname}{Inhaltsverzeichnis}


\newcommand{\absatz}{5mm}
\newcommand{\Blocksatz}{\justifying}
\newcommand{\strich}{\vspace{\absatz}\hrule\vspace{\absatz}} 


\title{\Thema}
\author{\authoren}
\date{\today}


\setlength{\parindent}{0cm} %deaktiviert die Automatische Einrückung am Anfang des Dokumentes 
\renewcommand{\contentsname}{\Inhaltsverzeichnisname}
\renewcommand{\dateseparator}{.}

%For subsubsub... sections
\setcounter{secnumdepth}{5}
\setcounter{tocdepth}{5}

\makeatletter
\newcommand\subsubsubsection{\@startsection{paragraph}{4}{\z@}{-2.5ex\@plus -1ex \@minus -.25ex}{1.25ex \@plus .25ex}{\normalfont\normalsize\bfseries}}
\newcommand\subsubsubsubsection{\@startsection{subparagraph}{5}{\z@}{-2.5ex\@plus -1ex \@minus -.25ex}{1.25ex \@plus .25ex}{\normalfont\normalsize\bfseries}}
\makeatother



\begin{document}
\begin{spacing}{\Zeilenabstand}
\maketitle
\begin{center}
    \vspace{-5mm}
    Fach: \Fach
\end{center}
\thispagestyle{empty}
\newpage

\pagestyle{fancy}
\fancyhead[L]{\authoren}
\fancyhead[C]{\Thema / Fach: \Fach}
\fancyhead[R]{\today}
\fancyfoot[R]{\thepage}
\fancyfoot[C]{}

\pagenumbering{roman}
\setcounter{page}{1}
\tableofcontents


\newpage

\pagenumbering{arabic}
\setcounter{page}{1}


%%% Document start %%%


\section{Git Herunterladen}

Der Git installer oder alternative installations möglichkeiten kann auf der \href{https://git-scm.com/install/windows}{Website} heruntergeladen werden

%\subsection{Windows}
%
%\subsubsection{Git Installer}
%
%Der Git Installer kann über die \href{https://git-scm.com/install/windows}{Website}
%heruntergeladen werden und ausgeführt werden.
%
%\subsubsection{Winget}
%
%Unter Windows kann mit der Shell(Cmd) und dem Programm Winget git insttaliert werden:
%
%\noindent 
%\begin{lstlisting}[language=bash]
%$ winget install git 
%\end{lstlisting}
%
%
%\subsection{MacOS}
%
%
%
%\subsection{Linux}

\section{Git}

\subsection{Git's Nutzen}

Git ist ein Programm zur Versionierung von Datein. Im ideal Fall sind diese Dateien
Textdateien git kann aber auch zur Versionierung von Binär Dateien genutz werden(Zum beispiel Bilder).

\subsection{Umgang mit Git}

\subsubsection{Git Repos}

Git funktioniert in sogenanten Reporsertories die man sich als git umgebung vorstellen kann.

\vspace{\absatz}

Mit der Cli kann man in einem Ordner den Comand {\color{red} git init} ein Git Reporsertorie erstellen
Im Hintergrund erstellt Git einen .git ordner in dem Git "lebt". Dieser Ordner ist in vielen 
Dateie Exploreren Versteckt und muss oftmals extra eingeblendet werden um ihn zu sehen.


\vspace{\absatz}


Diese Git Repo kann dann Beispielsweise über Github oder anderen Platformen online 
Syncronsiert werden. 

\vspace{\absatz}

Mit einem {\color{red}git clone} Klonnt man das Git repo vom Origin also z.B den stand von Github. Und 
macht eine Lokale Kopie. Werden dann änderungen auf diesem git repo gemacht sind diese änderungen
erstmalls lokal. Diese kann man dann mit push/pull/ und add Syncronsiert zum Origin(Github bspw.)






\subsubsection{Git Commands}

Git abreitet dann mit verschieden Operatoren:

\subsubsubsection{Clone}

\textbf{Clone:} 
Clone ist der befehl um ein git repo aus dem Internet zu Klonen. -> Auf das System zu schieben  




\subsubsubsection{Status}

Nach dem änderungen vorgenomen wurden kann man mit {\color{red}git status} diese änderungen einsehen.
Oder wenn in dem Git repo vom Origin verändert wurde. Also Commits hinter dem Aktuellen stand ist.

\subsubsubsection{Add}

\textbf{Add:}
Nachdem änderungen gemacht wurden können diese mit {\color{red}git add} zu einem Commit hinzugefügt werden, der 
syntax von dem Command ist wie folgt:

\vspace{\absatz}

Fügt eine Bestimmte Datei zum Commit hinzu

{\color{red}git add {filename}}

\vspace{\absatz}

Fügt alle geänderten Dateien zum Commit hinzu.

{\color{red}git add --all}

\subsubsubsection{Commit}

Git funktioniert in Commits also um zu den Git repo dazuzutragen muss eine Commit gemacht werden.
Erst müssen Dateien zu einem Commit mit {\color{red}git add} hinzugefügt werden um danach mit {\color{red} git commit -m ""}
zu Commiten. Das -m ist für die message/tittel des Commits


\subsubsubsection{Push}

Mit {\color{red}git push} kann man alle Commit die man erstellt hat auf den Origin vom Git repo Schieben
Syntax:

\subsubsubsection{Pull}

Mit Git Pull hollt man sich alle aktuellen Commits vom Origin die noch nicht auf dem Lokalen Git Repo 
liegen.














\subsubsection{Grafisch}






\end{spacing}
\end{document}